\documentclass{article}

\usepackage[preprint]{neurips_2026}
\usepackage[utf8]{inputenc}
\usepackage[T1]{fontenc}
\usepackage{booktabs}
\usepackage{microtype}
\usepackage{graphicx}
\usepackage{tabularx}
\usepackage{array}
\usepackage{url}
\usepackage{amsmath}
\usepackage{natbib}

\usepackage{comment}

%%%%%%%%%%%%%%%%%%%%%%%%%%%%%%%%%%%%%%%%%%%%%%%%%%%%%%%%%%%%
%%%%%%%%%%%%%%%%%%%%%%% Workshop %%%%%%%%%%%%%%%%%%%%%%%%
%%%%%%%%%%%%%%%%%%%%%%%%%%%%%%%%%%%%%%%%%%%%%%%%%%%%%%%%%%%%
% https://judge2026.github.io/
% - **Workshop:** JUDGe 2026 @ NeurIPS 2026
% - **Location:** Atlanta, Georgia
% - **Deadline:** August 29, 2026
% - **Full paper:** 6 pages + references
% - **Short paper:** 4 pages + references
% - **Review:** Double-blind via OpenReview
% - **Proceedings:** Non-archival
%%%%%%%%%%%%%%%%%%%%%%%%%%%%%%%%%%%%%%%%%%%%%%%%%%%%%%%%%%%%

%%%%%%%%%%%%%%%%%%%%%%%%%%%%%%%%%%%%%%%%%%%%%%%%%%%%%%%%%%%%
%%%%%%%%%%%%%%%%%%%%%%% Style Guide %%%%%%%%%%%%%%%%%%%%%%%%
%%%%%%%%%%%%%%%%%%%%%%%%%%%%%%%%%%%%%%%%%%%%%%%%%%%%%%%%%%%%
% - The official paper/venue style guide takes precedence over the rules below.
% - Use American English throughout.
% - Use \say{} for quotation marks (requires the dirtytalk package).
% - Include the GitHub link in the abstract.
% - Refer to figures as "Figure X", not "Fig. X".
% - Define abbreviations only once (e.g., SCM), and do not define abbreviations in the abstract.
% - Use a comma after "e.g." and "i.e." (e.g., blah blah blah; i.e., blah blah blah).
% - Use title case for paragraph, section, and subsection titles. See: https://capitalizemytitle.com/
% - End paragraph titles with a period.
% - Use \citep{} for "[Xu et al., 2026]" and \citet{} for "Xu et al. [2026]".
%%%%%%%%%%%%%%%%%%%%%%%%%%%%%%%%%%%%%%%%%%%%%%%%%%%%%%%%%%%%





\title{What Shapes the Verdict?\\Peculiarities of LLM-Based Persuasion Evaluation}

\author{
  Bhuvan Shingade\\
  RPTU Kaiserslautern-Landau\\
\texttt{bhuvan.shingade@edu.rptu.de}
  \And
  David A. Selby\\
  DFKI\\
\texttt{David.Selby@dfki.de}
  \And
  Sebastian Vollmer\\
  DFKI\\
\texttt{Sebastian.Vollmer@dfki.de}
  \And
  Gerrit Grossmann\\
  DFKI\\
\texttt{Gerrit.Grossmann@dfki.de} }

\begin{document}

\maketitle


\begin{abstract}
LLM-generated persuasion is increasingly relevant to democratic discourse, coordinated influence campaigns, and general attempts to shape public opinion. Because human evaluation is costly and difficult to scale, LLM judges are likely to be used to optimize and compare LLM-based persuasive strategies. We examine how reliably such judges rank debating models under different evaluation protocols. Using controlled changes to role assignment, speaking order, and candidate presentation order, we measure changes in model rankings, disagreement between judges, observational associations with response length, and the relative difficulty of defending each position. We find that rankings depend strongly on whether models defend opposing or matched positions, while position and candidate-order effects vary across judges. These results show that evaluation design materially shapes conclusions about model persuasiveness; they do not establish human persuasion effectiveness.
\end{abstract}

\section{Introduction}

LLMs can produce fluent, adaptive arguments in multi-turn settings. This makes them useful for benign applications such as explanation and correction, but also relevant to persuasion, misinformation, and coordinated influence. Human studies show that AI dialogue can change conspiracy beliefs and political preferences \citep{costello2024conspiracy,lin2025voters}. Recent large-scale experiments further indicate that prompting and post-training can increase conversational persuasiveness while reducing factual accuracy \citep{hackenburg2025levers}. Other work measures model persuasiveness directly \citep{durmus2024persuasiveness,singh2024persuasionbench} and analyzes how influence can propagate in multi-agent systems \citep{abedini2026stubborn}.

Because human evaluation is expensive, LLM judges are often used as scalable evaluators. Strong LLM judges can approximate human preferences in some pairwise comparison tasks \citep{zheng2023judge}. However, they are not neutral measurement devices. Prior work reports position bias, verbosity bias, self-preference, and other evaluator artifacts \citep{wang2024faireval,shi2024positionbias,panickssery2024selfpreference,chen2024judgementbias}. These issues are especially important for persuasion, where side framing, length, and rhetorical style may change the verdict.

We introduce a controlled benchmark that asks which design choices shape LLM-judged persuasive debates rather than treating the outcome as a final leaderboard. Rankings change between different-position and same-position comparison, and judge identity, position, speaking order, and candidate order remain consequential enough that aggregate win rates require explicit controls.


\section{Data Generation}


\subsection{Setup}

We use $X=10$ topics and $Y=4$ debating models: Apertus-70B, Llama-3.1-8B, Qwen-30B, and Gemma-31B. Each topic defines two symmetric propositions without assigning semantic meaning to their order. A debate contains eight alternating messages, four per model, limited to 220 tokens each.

\subsection{Debate Generation}
For one topic and two models, a debate condition is defined by the two propositions, the mapping from models to propositions, and the starting model. We generate all four combinations of the two possible model--proposition mappings and two possible starters for every unordered model pair. The design therefore contains $N=X\binom{Y}{2}\cdot4=10\cdot6\cdot4=240$ debates. Failed generations are repeated until every condition has one valid transcript.

\subsection{Judgment Generation}

Three judges evaluate the debates: Gemma-4-31B-IT, Qwen3-30B-A3B-Instruct-2507, and GLM-4.7. Because evaluator-specific behavior is central to the study, we report judges separately and use pooled summaries only where they clarify an overall pattern.

\paragraph{Mode 1: Different-Position Evaluation.}
Each judge selects which model argued its assigned position better in one transcript. This resembles an ordinary debate verdict, but confounds model quality with position difficulty. The mode yields $240\cdot3=720$ judgments.


\paragraph{Mode 2: Same-Position Evaluation.}
To reduce position confounding, we compare two role-swapped debates from the same topic and model pair. The same two models appear in both transcripts: each model defends the tested proposition once against the other model, under the same starting condition. The judge selects which candidate made the stronger case for that shared proposition, and every comparison is shown in both candidate orders. The resulting $10\binom{4}{2}\cdot2\cdot2\cdot2=480$ presentations represent 240 matched comparisons and yield $480\cdot3=1{,}440$ judgments.


\subsection{Protocol and Prompting Details}

\paragraph{Debate Protocol.}

A moderator presents the question, both propositions, and a neutral instruction to identify the central trade-off. The assigned starting model responds first, after which the models alternate for four turns each. Each model is instructed to maintain its assigned proposition, address the opponent's strongest recent point, and use concrete mechanisms, trade-offs, examples, or failure modes. References to judges, prompts, benchmarks, and model identities are prohibited.

\paragraph{Judging Protocol.}
Judges score argument quality, responsiveness, factual discipline, and persuasive clarity while ignoring personal policy preference. Same-position judges compare only the two candidates defending the specified proposition, and ties are allowed. Without external retrieval, factual discipline is assessed from the supplied topic and transcripts.

\section{Experimental Analysis}

\subsection{Experiment 1: Evaluation Design and Judge Agreement}

We ask whether rankings are stable across evaluation constructs. Different-position verdicts combine argument quality, position difficulty, and interaction dynamics; same-position comparisons hold the proposition and speaking condition fixed. Similar rankings would suggest that both modes recover a common notion of persuasive ability. Large rank changes would instead show that the benchmark conclusion depends on what the evaluation protocol treats as success.

\paragraph{Results.}
Table~\ref{tab:current-mode-ranks} reports rankings separately for each judge and mode. Within each judge and mode, models are ranked using the tie-adjusted score $(\mathrm{wins}+0.5\,\mathrm{ties})/\mathrm{evaluations}$; the final columns average ranks across judges. Rank 1 is best, and rows are ordered by mean different-position rank. The protocols do not produce a common ordering. Llama-3.1-8B has the best mean rank in different-position evaluation, although the Qwen judge ranks it last, but all three judges rank it last in same-position evaluation. Gemma-4-31B becomes strongest on average in the matched task, while Qwen3-30B remains comparatively stable in second place. Because same-position evaluation holds the defended proposition and starting condition fixed, this reversal shows that success in a complete debate does not transfer directly to the narrower ability to make the stronger case for a shared proposition. The protocols also differ in decisiveness: ties account for 22.8\% of different-position evaluations but only 0.3\% of same-position evaluations. The higher tie rate in complete debates is consistent with judges finding opposing-position performances harder to compare directly. A single pooled leaderboard would therefore conceal both a change in the measured construct and a large difference in judge decisiveness.

\begin{table}[t]
  \centering
  \footnotesize
  \setlength{\tabcolsep}{3pt}
  \caption{\textbf{Exp.~1 -- Model Ranking:} Per-judge ranks by evaluation mode (D: different-position; S: same-position).}
  \label{tab:current-mode-ranks}
  \begin{tabularx}{\textwidth}{@{}>{\raggedright\arraybackslash}Xcccccccc@{}}
    \toprule
    & \multicolumn{2}{c}{Gemma judge} & \multicolumn{2}{c}{Qwen judge}
    & \multicolumn{2}{c}{GLM judge} & \multicolumn{2}{c}{Mean rank} \\
    \cmidrule(lr){2-3}\cmidrule(lr){4-5}\cmidrule(lr){6-7}\cmidrule(l){8-9}
    Model & D & S & D & S & D & S & D & S \\
    \midrule
    Llama-3.1-8B & 1 & 4 & 4 & 4 & 1 & 4 & 2.00 & 4.00 \\
    Gemma-4-31B & 3 & 1 & 2 & 3 & 2 & 1 & 2.33 & 1.67 \\
    Qwen3-30B & 2 & 2 & 3 & 2 & 3 & 2 & 2.67 & 2.00 \\
    Apertus-70B & 4 & 3 & 1 & 1 & 4 & 3 & 3.00 & 2.33 \\
    \bottomrule
  \end{tabularx}
\end{table}

We next measure whether the verdict is robust to evaluator choice. Agreement is computed on identical comparison instances, so disagreement cannot be explained by different debates. High agreement would support treating judges as interchangeable evaluators; low agreement implies that a reported ranking is partly judge-specific.

\paragraph{Judge Agreement.}
All three judges give the same winner in 21 of 240 different-position rows (8.8\%) and 254 of 480 same-position rows (52.9\%). In the different-position mode, 48 rows have three distinct verdicts and therefore no majority winner; every same-position row has a majority. Thus, narrowing the comparison to a shared proposition makes judgments more consistent, but does not make the judges interchangeable: in nearly half of the matched presentations, at least one judge still disagrees. This is descriptive agreement among the three cloud judges, not agreement with human persuasion judgments.

\begin{comment}
\paragraph{Debater Confidence.}
Rate the confidence of each model in each debate and score it. Then see how well the judgment (who wins) is correlated with who sounded more confident. The current data do not contain a validated measure of how confident each debater sounds, so this relationship is not estimated here.

\paragraph{Confidence Diagnostic.}
The existing confidence fields measure judge confidence, not how confident the debaters sounded, and therefore cannot answer the proposed correlation. Mean judge confidence in different-position evaluation is 0.838 for Gemma, 0.709 for GLM, and 0.910 for Qwen; in same-position evaluation it is 0.967, 0.730, and 0.948, respectively. Without a human or otherwise validated correctness label, these values describe judge certainty but do not establish calibration.
\end{comment}

\begin{comment}
\paragraph{Repeated-Judgment Stability.}
This analysis requires additional judgments but no new debates. We re-evaluate all $N=240$ transcripts five times with each judge using the same judging prompt. We report within-judge agreement, the standard deviation of model win rates, and the frequency with which repeated judgments produce different winners.

\paragraph{Current Stability Diagnostic.}
The present files do not contain five repeated judgments with identical input order. They do permit a stricter candidate-order check using the same two underlying debates. The winning model remains unchanged in 190 of 240 paired reversals for Gemma (79.2\%), 91 for Qwen (37.9\%), and 166 for GLM (69.2\%). Because the manipulation changes candidate presentation order, this result is reported as an order effect in Experiment~2 rather than as stochastic repeatability.
\end{comment}

\begin{comment}
\paragraph{Sensitivity to Style.}
This analysis requires an additional perturbed version of the debate transcripts. For each transcript, we generate one style-normalized version that preserves the argument content while making both speakers use the same tone and formatting. We judge the original and modified transcripts and report the fraction of verdicts that change.

\paragraph{Current Status.}
No controlled style-normalized or paraphrased transcript variants are present in the curated files. Natural style differences are confounded with model identity and argument content, so style sensitivity cannot be estimated from the current data without the proposed perturbation experiment.
\end{comment}

\paragraph{Prompt-Phrasing Sensitivity.}
Judge instructions are part of the measurement instrument: semantically similar wording may change how a model weighs the same criteria. We therefore re-judge a balanced 12-debate subset with the canonical instruction and a meaning-preserving paraphrase, while holding transcripts, candidate labels, output format, temperature, and token limit fixed. Gemma and Qwen each agree across phrasings on 7 of 12 cases (58.3\%). Across 24 paired comparisons, the 10 changes comprise two decisive winner reversals and eight transitions between a winner and a tie. GLM could not complete the paired sample because of repeated API timeouts and is excluded. Because each prompt receives only one sample per case, this result is a small cloud diagnostic and cannot separate wording sensitivity from residual API nondeterminism.

\subsection{Experiment 2: Biases in LLM Persuasion Judging}

\paragraph{Position and Speaking-Order Bias.}
A persuasion benchmark should reward advocacy quality rather than an arbitrary proposition label or turn order. The factorial design crosses each model pair with both proposition assignments and both speaking orders. We report the larger of the two proposition win shares for each judge, which measures the magnitude of side imbalance without attaching meaning to either label. We also report the starting speaker's decisive win share. Because the starter does not deliver the final turn, this statistic contrasts opening with closing position and is not a pure recency estimate.

\paragraph{Results.}
Table~\ref{tab:current-position-order} shows no universal preference for one proposition label: the favored side differs across judges, and the magnitude of imbalance ranges from 52.5\% to 63.0\% of decisive verdicts. The pooled starting-speaker share is 45.5\%, with substantial judge variation and the strongest closing-speaker tendency under Qwen. These patterns indicate judge-specific sensitivity to side content and turn order rather than a general advantage for an arbitrarily named position. Candidate presentation order is also judge-specific (Table~\ref{tab:current-candidate-order}): Qwen reverses its selected model in most paired order swaps, whereas Gemma and GLM remain stable more often than not. The candidate-order comparison holds the underlying debates fixed and therefore directly diagnoses presentation-order sensitivity; comparisons across speaking-order cells can additionally reflect generation variation because those transcripts were generated separately.

\begin{table}[t]
  \centering
  \footnotesize
  \caption{\textbf{Exp.~2.1 -- Position and Speaking-Order Bias:} Favored-proposition and starting-speaker shares among decisive different-position verdicts. The favored proposition is defined separately for each row.}
  \label{tab:current-position-order}
  \begin{tabularx}{\textwidth}{@{}>{\raggedright\arraybackslash}Xrr@{}}
    \toprule
    Judge & Favored proposition (\%) & Starting speaker (\%) \\
    \midrule
    All judges & 54.5 & 45.5 \\
    Gemma-4-31B & 54.3 & 50.4 \\
    GLM-4.7 & 52.5 & 47.5 \\
    Qwen3-30B & 63.0 & 39.7 \\
    \bottomrule
  \end{tabularx}
\end{table}

\begin{table}[t]
  \centering
  \footnotesize
  \setlength{\tabcolsep}{4pt}
  \caption{\textbf{Exp.~2.1 -- Candidate-Order Stability:} Stability under candidate-order reversal in same-position evaluation. Each pair contains the same two underlying debates presented in both candidate orders.}
  \label{tab:current-candidate-order}
  \begin{tabularx}{\textwidth}{@{}>{\raggedright\arraybackslash}Xrrrrr@{}}
    \toprule
    Judge & Pairs & Stable & Reversals & Tie changes & Stable (\%) \\
    \midrule
    Gemma-4-31B & 240 & 190 & 46 & 4 & 79.2 \\
    GLM-4.7 & 240 & 166 & 74 & 0 & 69.2 \\
    Qwen3-30B & 240 & 91 & 149 & 0 & 37.9 \\
    \bottomrule
  \end{tabularx}
\end{table}

\paragraph{Verbosity Bias.}
Verbosity bias matters because a judge may mistake additional text for stronger reasoning, rewarding output length even when substantive quality is unchanged. For each decisive judgment, we compare the total word count produced by the selected model with that of the non-selected model. A larger mean among winners is the expected observational signature of verbosity preference, but it is not causal: stronger models or arguments may also be longer. A causal test requires paired variants that preserve argument content while equalizing length.

\paragraph{Results.}
Table~\ref{tab:current-verbosity} pools both evaluation modes and reports one row per judge. Winners are longer on average for all three judges, which is consistent with a verbosity preference but does not establish one: response length remains confounded with model identity, topic, and argument quality. A causal claim requires matched arguments whose content is preserved while length is varied. Values are word counts and the reported variances are sample variances in squared words.

\begin{table}[t]
  \centering
  \footnotesize
  \setlength{\tabcolsep}{4pt}
  \caption{\textbf{Exp.~2.2 -- Verbosity Bias:}\enspace Observed argument length for winners and losers, pooled across evaluation modes. This is an observational diagnostic.}
  \label{tab:current-verbosity}
  \begin{tabularx}{\textwidth}{@{}>{\raggedright\arraybackslash}Xrrrrr@{}}
    \toprule
    Judge & $n$ & Win mean & Win var. & Loss mean & Loss var. \\
    \midrule
    Gemma-4-31B & 603 & 639.0 & 2344.4 & 621.1 & 5284.3 \\
    GLM-4.7 & 720 & 635.3 & 2584.8 & 623.9 & 5202.9 \\
    Qwen3-30B & 669 & 639.4 & 2971.1 & 619.6 & 4831.4 \\
    \bottomrule
  \end{tabularx}
\end{table}

\paragraph{Self-Preference Bias.}
Self-preference would make evaluations systematically more favorable when the judge and a candidate share a model identity. For each comparison involving judge-model $X$, we measure how often $X$ selects itself and compare this with the selection rate of judges that generated neither candidate response on the same comparison. Using identical cases holds debate difficulty fixed, while excluding both participants prevents the reference group from containing a second self-judgment.

\paragraph{Results.}
Pooling both evaluation modes, Gemma selects itself in 60.3\% of relevant comparisons, compared with a 63.2\% selection rate from uninvolved judges on the same cases. Qwen selects itself in 52.2\% of comparisons, compared with a 65.4\% uninvolved-judge baseline. Neither model therefore shows positive self-preference under this diagnostic. The comparison remains descriptive because judges may apply different evaluation criteria.

\begin{comment}
\paragraph{Ideological Preference.}
This analysis requires an additional annotation of the debate propositions. We select the politically relevant topics and ask three independent annotator models to classify both propositions on a five-point scale from strongly liberal to strongly conservative. The majority annotation is used as the proposition label. We then compare, for each judge, the probability that the liberal or conservative proposition wins. Because each model argues both sides and both speaking orders, these effects can be estimated while controlling for debater identity and speaking order.

\paragraph{Current Status.}
The present topics do not have validated ideological annotations. The current data support judge-specific position preferences, but not a defensible claim about ideological preference. We therefore do not estimate this effect until an annotation protocol is fixed.
\end{comment}



\subsection{Experiment 3: Difficulty of Debate Positions}

Finally, we ask whether the substantive proposition being defended changes the chance of winning. This matters because a model leaderboard is misleading if a model succeeds mainly because it was assigned the more defensible side. For each topic, we pool verdicts across debating models, role assignments, speaking orders, and judges. We call the proposition with the larger decisive win share the \emph{observationally easier} position. Balancing model and order assignments makes this statistic less dependent on a particular debater, although judge preferences and generation variation still prevent a causal or human-validated interpretation.

\paragraph{Results.}
Table~\ref{tab:current-position-difficulty} shows the four largest observed departures from an even split. Energy infrastructure has the strongest observed asymmetry: decentralized renewable microgrids win 40 decisive evaluations, compared with 18 for centralized nuclear fission. This indicates that side difficulty can materially affect a model's apparent success even in a balanced design, although pooled judge verdicts do not establish which position humans would find easier. The current topic set also lacks a deliberately obvious control item, which should be added only after its wording and selection criteria are pre-specified.

\begin{table}[t]
  \centering
  \footnotesize
  \caption{\textbf{Exp.~3 -- Position Difficulty:} Topics with the largest observed asymmetries. Easy share is the easier proposition's share among decisive verdicts; the final column gives easy-position wins, hard-position wins, and ties.}
  \label{tab:current-position-difficulty}
  \begin{tabularx}{\textwidth}{@{}>{\raggedright\arraybackslash}p{0.17\textwidth}>{\raggedright\arraybackslash}X>{\raggedright\arraybackslash}Xrr@{}}
    \toprule
    Topic & Observationally easier & Observationally harder & Easy (\%) & E/H/T \\
    \midrule
    Energy infrastructure & Decentralized renewable microgrids & Centralized nuclear fission & 69.0 & 40/18/14 \\
    Mega-event hosting & Multi-country co-hosting & Single-nation hosting & 63.0 & 34/20/18 \\
    Cultural repatriation & Repatriate artifacts and use digital replicas & Retain artifacts in universal museums & 62.3 & 38/23/11 \\
    Exploration priority & Prioritize deep-sea exploration & Prioritize space exploration & 61.8 & 34/21/17 \\
    \bottomrule
  \end{tabularx}
\end{table}

\begin{comment}
\subsection{Experiment 4: Comparison with Human Debate Judgments}

As an external validation, we use the Debate.org (DDO) dataset (\url{https://esdurmus.github.io/ddo.html}), which contains human-written debates in which other human users voted on the winning side. More specifically, we use the cleaned subset introduced by Rescala et al.\ (EMNLP 2024) for studying whether LLMs can recognize convincing arguments. This subset contains 833 political debates with clear Pro/Con positions and retains the human votes indicating the winner. We apply the same LLM judges used in our experiments to these debate transcripts, without revealing the human votes, and compare their verdicts with the human majority winner. We report LLM--human agreement, inter-LLM agreement, and additionally test whether position, verbosity, and ideological biases observed in our generated debates also occur when judging human debates.
\end{comment}






\section{Related Work}


\paragraph{Persuasive and Strategic Dialogue.}
LLM persuasion is now studied both as a capability and as a risk. \citet{costello2024conspiracy} show that personalized GPT-4 Turbo dialogues reduced conspiracy belief by about 20\% on average in a human-subjects experiment. \citet{durmus2024persuasiveness} and \citet{singh2024persuasionbench} study the persuasiveness of model-generated arguments. Earlier negotiation work showed that dialogue agents can optimize strategic goals rather than merely imitate human language \citep{lewis2017deal}. Our setting differs because it evaluates model-vs-model policy debates using an automated judge rather than measuring human belief change.

Recent work connects persuasion directly to evaluator reliability. \citet{zahraei2026prior} show that LLM judges can conflate prior agreement with argumentative quality, including in debate evaluation. \citet{carro2025beliefs} compare sequential and simultaneous subjective debates and report a second-speaker effect in the sequential protocol. Multi-turn frameworks also study both persuasive effectiveness and susceptibility \citep{bozdag2025persuade}. Separately, \citet{hwang2025trick} demonstrate that rhetorical insertions can alter LLM-judge scores and pairwise verdicts, while \citet{bandaru2025political} use structured multi-agent debate to study how model and persona choices shape political attitudes. Together, this literature motivates treating the evaluation prompt, judge identity, and interaction protocol as parts of the measurement design.

\paragraph{Debate and Multi-Agent Influence.}
Debate has been proposed as a mechanism for eliciting useful information from competing agents \citep{irving2018debate}. Multi-agent work also shows that interaction structure can shape collective beliefs. For example, \citet{abedini2026stubborn} model LLM multi-agent belief propagation with Friedkin-Johnsen opinion dynamics and study cascade-like manipulation. Our benchmark is narrower: it uses two-agent debates, but it shares the concern that interaction design changes observed outcomes.

\paragraph{LLM-as-a-Judge Bias.}
\citet{zheng2023judge} introduce MT-Bench and Chatbot Arena and show that strong LLM judges can be useful for pairwise model comparison. Subsequent work highlights systematic limitations. \citet{wang2024faireval} study unfairness in LLM evaluation and propose position calibration. \citet{shi2024positionbias} analyze position bias across judges and tasks. \citet{panickssery2024selfpreference} show that LLM evaluators can favor their own generations, while \citet{chen2024judgementbias} compare human and LLM judging biases. These results motivate our controls for judge identity, role assignment, speaker order, and output parsing.

\section{Conclusions and Future Work}
This benchmark shows that an LLM-judged persuasion ranking is not a property of the debating model alone. Changing the comparison construct from a complete different-position debate to a matched same-position comparison substantially reorders the four models, and the three judges often disagree on identical inputs. Candidate presentation order is especially consequential for Qwen, while the prompt-phrasing pilot shows that even meaning-preserving judge instructions can move verdicts across the tie boundary or reverse a decisive winner. The observed association between response length and winning, the matched self-preference comparison, and the topic-level position asymmetries further illustrate why aggregate win rates require explicit controls.

These findings support a methodological conclusion rather than a claim about which model is most persuasive to people: evaluator identity, evaluation mode, position assignment, and presentation order are components of the measurement instrument. The current results are descriptive cloud-judge outcomes and have no human persuasion ground truth. A stronger validation should therefore repeat judgments under forced-choice and repeated-sampling protocols, add pre-specified topics with a clearly easier position, and compare model verdicts with human judgments.

The next benchmark iteration will also broaden the capability range by adding one substantially stronger and one substantially weaker model. The same six-model set will serve as both debaters and judges, enabling self-preference comparisons without changing the evaluator population across roles. Controlled length and style interventions, validated ideological annotations, and repeated prompt variants can then separate causal evaluator effects from the observational patterns reported here.





\bibliographystyle{plainnat}
\bibliography{references}

\end{document}





\begin{comment}


\begin{itemize}

    \item \textbf{Dataset I: Debate generation}
    \begin{itemize}
        \item Number of questions: $X$ (e.g., 20).
        \item Number of models: $Y$ (e.g., 4).
        \item Length of a debate: $Z$ rounds.
        \item Every debate is run four times:
        \begin{itemize}
            \item Model 1 starts vs.\ Model 2 starts.
            \item Model 1 argues for position A vs.\ Model 2 argues for position A.
        \end{itemize}
        \item Total number of debates:
        \[
            N = X \cdot \frac{Y(Y-1)}{2} \cdot 4.
        \]
        \item Example: $N = 480$.
        \item Debates that result in an error are re-run so that there is one valid debate for each condition.
        \item The debates are stored in \texttt{RawDebates.csv}.
        \item Columns: \texttt{DebateId}, \texttt{ModelAname}, \texttt{ModelBname},
        \texttt{positionA}, \texttt{positionB}, \texttt{DiscussionPlainText},
        and \texttt{starting}.
        \item Example row: \texttt{DebateId = 1}, \texttt{ModelAname = ChatGPT},
        \texttt{ModelBname = Grok}, \texttt{positionA = Dinosaurs are real},
        \texttt{positionB = Dinosaurs are fiction},
        \texttt{DiscussionPlainText = ModelA: Dinosaurs are supported by extensive fossil evidence.
        ModelB: Fossils could be interpreted as constructed or misleading evidence...},
        \texttt{starting = A}.
        \item \texttt{positionA} is fixed for a specific debate topic.
        \item A debate can start with position A or position B.
    \end{itemize}

    \item \textbf{Dataset II: Judge generation}
    \begin{itemize}
        \item Perform preference ranking using two evaluation modes.

        \item \textbf{Mode 1: Different-position evaluation}
        \begin{itemize}
            \item Evaluate who argued better: position A or position B.
            \item Store the results in \texttt{DiffPosJudgements.csv}.
            \item Columns: \texttt{DebateId}, \texttt{Judge 1 (Grok) winner},
            \texttt{Judge 2 (ChatGPT) winner}, and
            \texttt{Judge 3 (Gemini) winner}.
            \item Example row: \texttt{DebateId = 1},
            \texttt{Judge 1 (Grok) winner = ChatGPT},
            \texttt{Judge 2 (ChatGPT) winner = Grok},
            \texttt{Judge 3 (Gemini) winner = ChatGPT}.
            \item Additional judge conditions can be included, e.g.,
            factualness vs.\ rhetoric.
            \item The file contains $N$ rows, with
            $N \times \texttt{JudgeNumber}$ evaluations.
        \end{itemize}

        \item \textbf{Mode 2: Same-position evaluation}
        \begin{itemize}
            \item Evaluate who argued better for the same position.
            \item Store the results in \texttt{SamePosJudgements.csv}.
            \item Columns: \texttt{DebateId1}, \texttt{DebateId2},
            \texttt{positionA}, \texttt{positionB}, \texttt{testedPosition},
            \texttt{Judge 1 (Grok) winner}, \texttt{Judge 2 (ChatGPT) winner},
            and \texttt{Judge 3 (Gemini) winner}.
            \item Example row: \texttt{DebateId1 = 2}, \texttt{DebateId2 = 3},
            \texttt{positionA = Dinosaurs are real},
            \texttt{positionB = Dinosaurs are fiction},
            \texttt{testedPosition = Dinosaurs are fiction},
            \texttt{Judge 1 (Grok) winner = ChatGPT},
            \texttt{Judge 2 (ChatGPT) winner = Grok},
            \texttt{Judge 3 (Gemini) winner = Grok}.
        \end{itemize}

    \end{itemize}

\end{itemize}



\section{Experimental Analysis}

\begin{itemize}

    \item \textbf{Experiment 1: Evaluation design and judge agreement}
    \begin{itemize}

        \item \textbf{Comparison of evaluation protocols}
        \begin{itemize}
            \item Opposing-position protocol: two models argue for different sides
            and the judge selects the better argument.
            \item Same-position protocol: two models separately defend the same
            position against a shared opponent and the judge compares their performance.
            \item \textbf{Tested:} Whether the two protocols produce similar model
            rankings and pairwise preferences.
            \item \textbf{Reported:} Rank correlation, pairwise ranking reversals,
            win-rate differences, and the fraction of model pairs for which the
            preferred model changes between protocols.
        \end{itemize}

        \item \textbf{Agreement between judge models}
        \begin{itemize}
            \item Evaluate every debate using several LLM judges while keeping the
            debating models, prompts, and debate transcripts fixed.
            \item \textbf{Tested:} Whether different judge models select the same
            winners and produce similar overall rankings.
            \item \textbf{Reported:} Pairwise judge agreement, majority agreement,
            rank correlation between judges, and the fraction of verdicts that
            depend on the selected judge.
        \end{itemize}

        \item \textbf{Stability across repeated judgments}
        \begin{itemize}
            \item Ask each judge to evaluate the same transcript multiple times
            using identical or minimally varied evaluation prompts.
            \item \textbf{Tested:} Whether judgments are stable across repeated
            samples and small changes to the judging prompt.
            \item \textbf{Reported:} Within-judge agreement, variance in model
            win rates, and changes in rankings across repeated evaluations.
        \end{itemize}

        \item \textbf{Sensitivity to style}
        \begin{itemize}
            \item Test how sensitive the judgment is to small changes in the text,
            e.g., changing the style slightly.
            \item Ideally, stronger models are less sensitive to these changes.
        \end{itemize}

    \end{itemize}


    \item \textbf{Experiment 2: Biases in LLM persuasion judging}
    \begin{itemize}

        \item \textbf{Position and speaking-order bias}
        \begin{itemize}
            \item Repeat each debate after swapping the positions assigned to the models.
            \item Repeat each debate after reversing which model speaks first or
            gives the final statement.
            \item \textbf{Tested:} Whether a model benefits from defending a
            particular position, speaking first, or speaking last.
            \item \textbf{Reported:} Win-rate changes after each reversal,
            verdict-flip rates, and the estimated advantage associated with each
            position and speaking slot.
        \end{itemize}

        \item \textbf{Verbosity bias}
        \begin{itemize}
            \item Measure the relationship between argument length and judge
            preference in unconstrained debates.
            \item Repeat a subset of debates with matched length limits or with
            shortened versions of the same arguments.
            \item \textbf{Tested:} Whether judges systematically prefer longer
            answers independently of their argumentative quality.
            \item \textbf{Reported:} The association between response length and
            winning, win-rate differences under matched-length conditions, and
            verdict changes caused by shortening an argument.
        \end{itemize}

        \item \textbf{Ideological preference of judges}
        \begin{itemize}
            \item Use politically or ideologically sensitive propositions.
            \item Repeatedly swap which debating model defends each side.
            \item \textbf{Tested:} Whether a judge systematically prefers particular
            political or ideological positions after controlling for the identity
            of the debating model, speaking order, and debate role.
            \item \textbf{Reported:} Position-specific win rates for each judge,
            estimated ideological preference effects, verdict-flip rates under
            role reversal, and differences in these effects between judge models.
        \end{itemize}

        \item \textbf{Self-preference bias}
        \begin{itemize}
            \item Test the same effects for self-preference bias.
        \end{itemize}

    \end{itemize}


    \item \textbf{Experiment 3: Difficulty of debate positions}
    \begin{itemize}
        \item Test whether one side of a proposition is consistently easier to
        argue for than the other.
        \item Reuse the debates from the previous experiments.
        \item Estimate the effect of defending each position while accounting for
        the debating model, judge model, speaking order, response length, and
        evaluation protocol.
        \item \textbf{Tested:} Whether particular positions achieve systematically
        higher win rates after controlling for known judge biases and experimental
        design choices.
        \item The analysis closely follows the ideological-bias analysis but asks
        whether the preference is consistent across judges rather than specific
        to a judge model.
        \item \textbf{Reported:} Adjusted win rates for both sides of each proposition,
        the estimated advantage of each position, uncertainty intervals, and the
        consistency of the effect across debating models, judges, and evaluation
        protocols.
    \end{itemize}

\end{itemize}



\section{Method}

\subsection{Debate Task}

Each benchmark case contains a question and two symmetric sides, called \emph{Position A} and \emph{Position B}. Agent A defends Position A, and Agent B defends Position B. We avoid a simple for/against framing because many policy choices compare two positive alternatives rather than an affirmative claim and its negation. For example, one energy topic compares centralized nuclear fission with decentralized renewable microgrids.

Debates begin with a moderator opening, followed by alternating agent turns and closing statements. The reported experiments use short debates: usually two turns and 60--140 maximum tokens per response. This keeps cost low and makes the experiments diagnostic rather than a final test of long-form deliberation.

\subsection{Models and Judges}

The benchmark supports local Ollama models and OpenAI-compatible cloud models. Local pilots use small open-weight models, including Qwen-0.5B, TinyLlama-1.1B, Gemma-2B, Llama-3B, and Qwen-3B. Cloud experiments use Academic Cloud/SAIA models. The main cloud judge is Gemma-31B; the main cloud debaters in the mode comparison are Apertus-70B, Llama-3.1-8B, and Qwen-30B.

We separate evidence by judge source. Cloud-judged results are treated as the strongest current evidence. Local Ollama judge results are diagnostic unless reproduced with a cloud judge.

\subsection{Debate Modes}

The benchmark implements three relevant comparison modes. In \emph{single} mode, one model-role assignment is evaluated. In \emph{paired} mode, the model-role assignment is reversed, so each model argues both positions. In \emph{same-position} mode, two candidate models each argue the same target position against a shared fixed opponent in separate debates, and the judge compares which candidate defended that target position better.

Same-position mode is designed to reduce confounding between model quality and side difficulty. It also requires careful interpretation, because the candidate comparison labels are distinct from the original debate-side labels.

\subsection{Controls and Evaluation Protocols}

Speaker order is controlled with two orders: Agent A first or Agent B first. The balanced setting runs both. This matters because a model may benefit from opening, closing, or appearing in a particular transcript slot.

The judge uses one of four protocols: holistic persuasion, argument quality, evidence/fact-checking, and deliberative quality. The same transcript can be judged under all four protocols. For this reason, we report these rows as judge evaluations rather than independent debates.







\section{Experiments}

Table~\ref{tab:experiment-overview} summarizes the main runs. We distinguish cloud-judged runs, which provide the strongest current evidence, from local diagnostic pilots, which are used to identify failure modes in the evaluation pipeline.

\begin{table}[t]
  \centering
  \small
  \caption{Overview of the main experiments. ``Eval.'' denotes judge evaluations after protocol expansion.}
  \label{tab:experiment-overview}
  \begin{tabular}{llrrl}
    \toprule
    Experiment & Judge & Eval. & Valid & Main purpose \\
    \midrule
    Local baseline ladder & Qwen-3B & 144 & 126 & Local sanity check \\
    Cloud baseline ladder & Gemma-31B & 144 & 144 & Parseability and weak baseline \\
    Mode comparison & Gemma-31B & 84 & 80 & Opposing vs same-position \\
    Position/order diagnostic & Gemma-31B & 48 & 48 & Side and speaker-order effects \\
    Local diagnostic pilots & Local judges & 96 & 89 & Judge fragility and bias probes \\
    \bottomrule
  \end{tabular}
\end{table}

\subsection{Baseline Capacity Ladder}

The baseline ladder compares adjacent local models on three policy topics: political-ad verification, de-extinction ethics, and zero-carbon energy infrastructure. The local pilot used Qwen-3B as judge, paired role swaps, balanced speaker order, all four protocols, two turns, and a 70-token limit.

The local judge produced 126 valid evaluations out of 144. It also showed severe artifacts: 18 malformed verdicts and 92 Agent A wins among valid rows. The same local run also showed a length signal, with the longer response winning 80 of 126 comparable valid evaluations. We therefore treat this run as a validity diagnostic rather than a model ranking.

The same ladder was rerun with the cloud Gemma-31B judge. This produced 144 valid verdicts out of 144. Position A and Position B each won 37 evaluations, with 70 ties. Thus, the cloud judge removed the local parse failures and did not reproduce the strong Agent A artifact. Table~\ref{tab:cloud-baseline} summarizes the model-level results.

\begin{table}[t]
  \centering
  \small
  \caption{Cloud-judged baseline ladder. Rows are judge evaluations; ties are included separately.}
  \label{tab:cloud-baseline}
  \begin{tabular}{lrrrrr}
    \toprule
    Model & Eval. & Wins & Losses & Ties & Win Rate \\
    \midrule
    Gemma-2B & 96 & 53 & 7 & 36 & 0.55 \\
    Qwen-0.5B & 48 & 14 & 0 & 34 & 0.29 \\
    Llama-3B & 48 & 7 & 5 & 36 & 0.15 \\
    TinyLlama-1.1B & 96 & 0 & 62 & 34 & 0.00 \\
    \bottomrule
  \end{tabular}
\end{table}

TinyLlama-1.1B lost every decisive cloud-judged row and is a reliable weak baseline in this setup. Gemma-2B ranks highest in the aggregate ladder, mainly because it decisively beats TinyLlama-1.1B. The Gemma-2B versus Llama-3B pair remains unresolved: Llama-3B leads 7--5, but 36 of 48 evaluations are ties. The high tie rate suggests that the 70-token limit may be too restrictive for fine-grained model separation.

Pair-level interpretation is more reliable than the aggregate ladder because the design uses adjacent comparisons rather than all-vs-all matches. In the cloud run, Qwen-0.5B beats TinyLlama-1.1B by 14--0 with 34 ties; Gemma-2B beats TinyLlama-1.1B by 48--0 with no ties; and Llama-3B narrowly leads Gemma-2B by 7--5 with 36 ties. Thus, only the TinyLlama comparisons produce a clean separation.

\subsection{Opposing-Position vs Same-Position Comparison}

The main cloud comparison evaluates Apertus-70B against Llama-3.1-8B on the same three policy topics. The opposing-position condition uses paired role swaps. The same-position condition uses Qwen-30B as a fixed opponent, while both candidates separately defend the same target side. All runs use the Gemma-31B judge, balanced speaker order, four protocols, two turns, and a 120-token limit.

The experiment produced 80 valid evaluations and 4 failed rows due to cloud rate limits. Apertus-70B won 65 of 80 valid evaluations. Table~\ref{tab:mode-comparison} shows the mode split.

\begin{table}[t]
  \centering
  \small
  \caption{Cloud comparison of opposing-position and same-position designs.}
  \label{tab:mode-comparison}
  \begin{tabular}{lrrp{0.32\linewidth}}
    \toprule
    Mode & Apertus & Llama & Note \\
    \midrule
    Opposing-position & 35 & 13 & Complete three-topic run \\
    Same-position, target B & 23 & 1 & Complete three-topic run \\
    Same-position, target A & 7 & 1 & Partial; rate-limited after one topic \\
    \midrule
    Total & 65 & 15 & 80 valid evaluations \\
    \bottomrule
  \end{tabular}
\end{table}

Apertus-70B ranks above Llama-3.1-8B in both designs and under all four protocols. Same-position comparison gives a sharper separation, especially when both candidates defend Position B. However, the Position A same-position condition is partial, so we do not claim full protocol agreement, rank correlation, or reversal statistics.

The protocol-level ranking is stable in the completed rows: Apertus-70B is preferred over Llama-3.1-8B under holistic persuasion, argument quality, evidence/fact-checking, and deliberative quality. The main technical difference is therefore not a rank reversal, but the margin: the opposing-position design gives Apertus 35 of 48 wins, while the completed same-position Position B condition gives Apertus 23 of 24 wins.

\subsection{Position and Speaking-Order Effects}

The cloud paired run also tests side and order effects. Across 48 valid evaluations, Position B wins 33 and Position A wins 15. The split by speaker order is also uneven: when Agent A speaks first, Agent B wins 22 of 24 evaluations; when Agent B speaks first, the split is closer.

\begin{table}[t]
  \centering
  \small
  \caption{Position and speaking-order diagnostic from the cloud paired run.}
  \label{tab:position-order}
  \begin{tabular}{lrr}
    \toprule
    Diagnostic & Count & Total \\
    \midrule
    Position A wins & 15 & 48 \\
    Position B wins & 33 & 48 \\
    Agent A wins when A speaks first & 2 & 24 \\
    Agent B wins when A speaks first & 22 & 24 \\
    Agent A wins when B speaks first & 13 & 24 \\
    Agent B wins when B speaks first & 11 & 24 \\
    Longer response wins & 21 & 48 \\
    \bottomrule
  \end{tabular}
\end{table}

This does not prove a universal Position B advantage. It does show that side and order effects can be large even when the judge is parseable and the design is balanced. Aggregate model win rates should therefore be reported with side- and order-specific diagnostics.

The length diagnostic in the same run is comparatively weaker: the longer response wins 21 of 48 evaluations. This suggests that, in this cloud run, side and order effects are more salient than a simple verbosity effect.

\subsection{Position Difficulty}

The same 48 cloud evaluations provide a descriptive position-difficulty diagnostic. All three topics favor Position B in this run, although the evidence/fact-checking protocol favors Position A on the energy topic.

\begin{table}[t]
  \centering
  \small
  \caption{Descriptive position difficulty in the cloud diagnostic.}
  \label{tab:position-difficulty}
  \begin{tabular}{lrrp{0.30\linewidth}}
    \toprule
    Topic & A Wins & B Wins & Pattern \\
    \midrule
    Political ads & 4 & 12 & Strong Position B advantage \\
    De-extinction & 5 & 11 & Position B advantage \\
    Energy grid & 6 & 10 & Position B advantage; protocol variation \\
    \bottomrule
  \end{tabular}
\end{table}

These are descriptive counts, not regression-adjusted causal estimates. They nevertheless show why topic-side difficulty should be reported rather than averaged away.

The protocol split is informative. Political-ad verification favors Position B under all four protocols. Energy infrastructure is less uniform: evidence/fact-checking favors Position A 3--1, while the other protocols favor Position B 3--1. This indicates that evaluation protocol can affect fine-grained topic conclusions even when the aggregate side pattern is stable.

\subsection{Local Diagnostic Pilots}

We also ran local pilots for judge agreement, repeated judgment, style sensitivity, verbosity, ideology, and self-preference. These are not paper-grade findings about model quality. They are stress tests for the evaluation pipeline.

\begin{table}[t]
  \centering
  \small
  \caption{Local diagnostic pilots. These runs are used to identify artifacts, not to make final model claims.}
  \label{tab:local-pilots}
  \begin{tabular}{lrrl}
    \toprule
    Pilot & Eval. & Valid & Main observation \\
    \midrule
    Judge agreement & 36 & 29 & 26 Agent A wins; 7 malformed rows \\
    Repeated judgment & 12 & 12 & 4 of 4 transcript groups stable \\
    Style sensitivity & 12 & 12 & 12 Agent A wins \\
    Verbosity & 16 & 16 & Longer response wins 7 of 16 \\
    Ideology/self-preference & 32 & 32 & 32 Agent A wins \\
    \bottomrule
  \end{tabular}
\end{table}

The pattern is clear: local judges can be stable on a tiny repeated sample, but they can also collapse into label preference. The style and ideology/self-preference pilots are dominated by Agent A wins, so they cannot support claims about style robustness, ideological preference, or self-preference. The verbosity pilot gives no strong length effect, but it only varies token limits and does not shorten the same argument post hoc.

\end{comment}
